\documentclass[12pt]{article}

\usepackage[a4paper, margin=1in]{geometry}

\begin{document}

% ---------- COVER PAGE ----------
\begin{titlepage}

\centering

\vspace*{2cm}

{\Huge \textbf{TECHNICAL UNIVERSITY OF}}\\[0.2cm]
{\Huge \textbf{KOSICE}}

\vspace{6cm}

{\LARGE \textbf{Habit Tracker Application}}

\vfill

\begin{flushleft}

\textbf{Name:} Sri Vimal Sanmugam \\[0.3cm]

\textbf{Subject:} Programming \\[0.3cm]

\textbf{Supervisor:} Marek Horvath

\end{flushleft}

\vspace{1cm}

2026

\end{titlepage}

\tab

\tableofcontents
\newpage

\section{Abstract}

This project describes the development of a Habit Tracking API using the Django web framework. 
The system allows users to create, manage, and track habits using a REST-style API. 
Instead of using a traditional database, the project stores all data in a JSON file 
called \texttt{storage.json}. 

The project demonstrates important backend development concepts such as CRUD operations, 
persistent storage, error handling, input validation, sorting, searching, exporting data, 
and resetting stored information. 

The API was tested using Postman and implemented through a sequence of tasks designed 
to gradually build a complete backend system.

\section{Introduction}

Habit tracking applications help users build productive routines by recording daily activities 
such as exercising, reading, or studying. Many modern applications rely on backend APIs to manage 
such data. APIs allow different applications to communicate using HTTP requests and JSON responses.

In this project, a Habit Tracking API was developed using Django. The API supports operations such 
as creating habits, retrieving stored habits, updating habit information, and deleting habits.

Instead of storing data in a database system such as MySQL or PostgreSQL, this project uses 
file-based persistence through a JSON file. This approach simplifies storage and allows the 
system to maintain data even after restarting the server.

\section{Project Objectives}

The main objectives of this project were:

\begin{itemize}

\item Understand how REST APIs operate
\item Learn backend development using Django
\item Implement CRUD operations (Create, Read, Update, Delete)
\item Implement persistent data storage using JSON files
\item Add input validation and error handling
\item Implement search and sorting features
\item Export stored data
\item Reset the stored data

\end{itemize}

The project also helps in understanding how backend systems manage data and handle client requests.

\section{Technologies Used}

\subsection{Python}

Python was used as the primary programming language due to its simplicity and wide usage 
in backend development.

\subsection{Django}

Django is a high-level Python web framework that allows developers to build secure 
and maintainable web applications quickly. Django provides built-in tools for handling 
requests, routing URLs, and generating responses.

\subsection{JSON}

JSON (JavaScript Object Notation) was used as the data storage format. 
All habit data is stored in the file \texttt{storage.json}.

\subsection{Postman}

Postman was used to test API endpoints. It allows sending HTTP requests such as GET, 
POST, PUT, and DELETE and inspecting server responses.

\subsection{Visual Studio Code}

Visual Studio Code was used as the development environment for writing and managing 
the project code.

\section{System Architecture}

The architecture of the system is relatively simple and consists of three main components:

\begin{itemize}

\item Client (Postman or Web Browser)
\item Django Backend Server
\item JSON File Storage

\end{itemize}

The workflow is as follows:

\begin{enumerate}

\item The client sends an HTTP request to the API.
\item Django processes the request.
\item The server reads or updates data in the JSON file.
\item The server sends a JSON response back to the client.

\end{enumerate}

\section{API Endpoints}

The following API endpoints were implemented in the project.

\section{API Endpoints}

The following API endpoints were implemented in the project.

\begin{center}
\begin{tabular}{|l|l|l|}
\hline
\textbf{Method} & \textbf{Endpoint} & \textbf{Description} \\
\hline
POST & /api/habit/ & Create a new habit \\
GET & /api/habit/all/ & Retrieve all habits \\
GET & /api/habit/\{id\}/ & Retrieve habit by ID \\
PUT & /api/habit/update/\{id\}/ & Update habit \\
DELETE & /api/habit/delete/\{id\}/ & Delete habit \\
GET & /api/habit/count/ & Return number of habits \\
GET & /api/habit/export/ & Download JSON data \\
DELETE & /api/habit/reset/ & Clear storage \\
\hline
\end{tabular}
\end{center}

\section{Project Implementation}

The project was developed through a sequence of tasks.

\subsection{Task 1: Initial API Creation}

The first step involved creating a Django project and defining basic API endpoints. 
Habit objects were initially stored in a Python list.

\subsection{Task 2: Loading Data on Startup}

A function was implemented to load data from the JSON file when the server starts. 
This allows the system to restore previously saved data.

\subsection{Task 3: Saving Data to File}

Whenever a new habit is created, the updated list of habits is written to the JSON file.

\subsection{Task 4: Updating Objects}

An endpoint was implemented to modify existing habit objects.

\subsection{Task 5: Deleting Objects}

An endpoint was created to delete habits and update the JSON file accordingly.

\subsection{Task 6: Safe File Loading}

Error handling was added to ensure the application does not crash if the JSON file does not exist.

\subsection{Task 7: Automatic ID Generation}

Unique IDs were automatically generated using the maximum existing ID plus one.

\subsection{Task 8: Retrieve Object by ID}

An endpoint was implemented that retrieves a habit using its ID.

\subsection{Task 9: Input Validation}

Validation checks were added to ensure required fields are provided.

\subsection{Task 10: Handling Invalid JSON}

Exception handling was implemented to handle invalid JSON input.

\subsection{Task 11: Search Function}

Users can search habits using query parameters.

Example:

\begin{verbatim}
GET /api/habit/all/?name=run
\end{verbatim}

\subsection{Task 12: Object Count}

An endpoint returns the total number of stored habits.

\subsection{Task 13: Sorting}

Habits can be sorted using query parameters such as:

\begin{verbatim}
?sort=id
?sort=-id
\end{verbatim}

\subsection{Task 14: Export Data}

Users can download the JSON file containing all stored habits.

\subsection{Task 15: Reset Storage}

An endpoint clears all stored data and resets the JSON file.

\section{Challenges Faced}

During development several issues occurred:

\begin{itemize}

\item Internal server errors caused by undefined variables
\item JSON decoding errors when invalid JSON was sent
\item Issues with ID generation after deleting objects
\item Understanding query parameters for sorting and search
\item Debugging API errors in Postman

\end{itemize}

These issues were resolved through debugging and proper exception handling.

\section{Testing}

All endpoints were tested using Postman.

The following tests were performed:

\begin{itemize}

\item Creating new habits
\item Retrieving stored habits
\item Updating habit information
\item Deleting habits
\item Searching habits
\item Sorting habits
\item Exporting JSON data
\item Resetting storage

\end{itemize}

All operations produced correct results.

\section{Results}

The final API successfully performs all required operations while maintaining 
persistent storage in a JSON file.

The system supports full CRUD operations and additional features such as search, 
sorting, exporting, and resetting data.

\section{Security Note}

In this project, encryption was implemented to protect sensitive data stored in files. 
The file \texttt{storage.json}, which contains user habit data, was encrypted using 
Fernet symmetric encryption from the cryptography library.

The encryption key is generated and stored in a separate file named 
\texttt{secret.key} during development. This ensures that the key is not mixed 
directly with the application logic.

To prevent exposing sensitive information, the key file is excluded from version 
control using a \texttt{.gitignore} file. This avoids accidental leakage of the 
encryption key when the project is uploaded to Git repositories.

This implementation is intended for educational purposes only. In real-world 
applications, more secure key management techniques would be used, such as 
environment variables or dedicated secret management systems.


\section{Authentication Implementation}

\subsection{Introduction}
In Week 6, authentication was introduced into the Habit Tracking API. The goal was to ensure that only authenticated users can access certain endpoints and that user-specific data is protected.

\subsection{Login Endpoint}
A login endpoint was implemented to allow users to authenticate using their username and password. Upon successful authentication, a token is generated and returned.

\textbf{Endpoint:}
\begin{verbatim}
/login/
\end{verbatim}

\textbf{Method:} POST

\textbf{Request Body:}
\begin{verbatim}
{
"username": "admin",
"password": "1234"
}
\end{verbatim}

\textbf{Successful Response:}
\begin{verbatim}
{
"message": "Login successful",
"token": "generated_token_here"
}
\end{verbatim}

\subsection{Token-Based Authentication}
A custom token model was used to store authentication tokens associated with users. This token is required to access protected endpoints.

\textbf{Header Format:}
\begin{verbatim}
Authorization: <token>
\end{verbatim}

\subsection{Protected Endpoint}
The \texttt{/habits/} endpoint was protected so that only authenticated users can access their own habits.

\textbf{Validation Steps:}
\begin{itemize}
\item Check if the Authorization token is present
\item Validate the token using the database
\item Retrieve the corresponding user
\item Filter habits belonging only to that user
\end{itemize}

\subsection{Example Code Snippet}
\begin{verbatim}
token = request.META.get("HTTP_AUTHORIZATION")

if not token:
return JsonResponse({"error": "Token required"}, status=401)

try:
token_obj = AuthToken.objects.get(token=token)
except AuthToken.DoesNotExist:
return JsonResponse({"error": "Invalid token"}, status=401)

user = token_obj.user
\end{verbatim}

\subsection{Error Handling}
The following authentication-related errors were handled:
\begin{itemize}
\item Missing token (401 Unauthorized)
\item Invalid token (401 Unauthorized)
\item Invalid login credentials (401 Unauthorized)
\item Missing username or password (400 Bad Request)
\item Invalid JSON format (400 Bad Request)
\item Incorrect HTTP method (405 Method Not Allowed)
\end{itemize}

\subsection{Testing}
Authentication was tested using Postman with the following scenarios:
\begin{itemize}
\item Successful login with valid credentials
\item Failed login with incorrect password
\item Missing fields in login request
\item Accessing protected endpoints without token
\item Accessing protected endpoints with valid token
\end{itemize}

\section{Encryption and Data Protection}

This week focused on improving the security of the application by introducing encryption and secure handling of sensitive data.

\subsection{Task 1: Installing Cryptography Library}

The \texttt{cryptography} library was installed using pip. The installation was verified by successfully importing the Fernet module, which is used for symmetric encryption.

\subsection{Task 2: Encryption Utility Implementation}

A separate Python file named \texttt{crypto\_utils.py} was created. This file contains functions for:

\begin{itemize}
\item Generating an encryption key
\item Loading the encryption key from a file
\item Encrypting files
\item Decrypting files
\end{itemize}

This modular approach improves code readability and maintainability.

\subsection{Task 3: File Encryption Testing}

A sample JSON file (\texttt{storage.json}) was created and encrypted into a binary file (\texttt{storage\_encrypted.bin}). The encrypted file was not readable in plain text, confirming that encryption was successful. The file was then decrypted back to its original form to verify correctness.

\subsection{Task 4: Encrypted Data Export}

The existing export functionality was extended by adding a new endpoint:

\begin{verbatim}
/api/habit/export-encrypted/
\end{verbatim}

This endpoint encrypts the stored data before exporting it. As a result, sensitive data is protected and cannot be read directly without decryption.

\subsection{Task 5: Handling Incorrect Keys}

Error handling was implemented to manage incorrect decryption keys. When a wrong key is provided, the system does not crash. Instead, it catches the exception and returns a clear error message:

\begin{verbatim}
Invalid key. Cannot decrypt data.
\end{verbatim}

This ensures predictable and secure behavior.

\subsection{Task 6: Protecting Sensitive Files}

The encryption key file (\texttt{secret.key}) was excluded from version control by adding it to the \texttt{.gitignore} file. This prevents accidental exposure of sensitive information when the project is shared.

\subsection{Task 7: Security Considerations}

A security note was added to explain the encryption approach used in the project. The note highlights that:

\begin{itemize}
\item Sensitive data is encrypted using Fernet symmetric encryption
\item The key is stored separately from the main code
\item The key is not included in version control
\item This implementation is intended for educational purposes
\end{itemize}

In real-world applications, more advanced key management techniques such as environment variables or secret management systems would be used.

\section{Automated Testing and Validation}

In Week 8, automated testing was implemented for the Habit Tracking System using Django's built-in \texttt{TestCase} framework. The goal was to ensure correctness, reliability, and robustness of the application through systematic testing.

\subsection{Overview}
A total of 12 test cases were implemented covering authentication, CRUD operations, filtering, file persistence, and system stability. All tests were executed successfully.

---

\subsection{Task 1: Create Tests Module}
A dedicated \texttt{tests/} folder was created to organize test cases. The tests were divided into:
\begin{itemize}
    \item \texttt{test\_auth.py} for authentication testing
    \item \texttt{test\_habit.py} for habit-related functionality
\end{itemize}
This modular structure improves readability and maintainability.

---

\subsection{Task 2: Basic Endpoint Testing}
Basic tests were written to verify that API endpoints respond correctly. These tests checked:
\begin{itemize}
    \item API accessibility
    \item Correct HTTP status codes (e.g., 200 OK)
\end{itemize}

---

\subsection{Task 3: Create Habit Testing}
Tests were implemented for habit creation functionality:
\begin{itemize}
    \item Valid input data returns success (status 200/201)
    \item Invalid or missing input fields return error (status 400)
\end{itemize}
This ensures proper validation and prevents invalid data from being stored.

---

\subsection{Task 4: Search and Filtering Testing}
Search and filtering features were tested by creating multiple habits with different attributes. The following were verified:
\begin{itemize}
    \item Search by habit name
    \item Filter by completion status (\texttt{is\_completed})
    \item Filter by frequency
\end{itemize}
These tests confirm that the system returns accurate filtered results.

---

\subsection{Task 5: Authentication Testing}
Token-based authentication was tested with the following scenarios:
\begin{itemize}
    \item Successful login returns a valid token (status 200)
    \item Invalid login returns unauthorized response (status 401)
    \item Requests without a token are rejected (status 401)
\end{itemize}
This ensures secure access to protected endpoints.

---

\subsection{Task 6: File Persistence Testing}
The system uses file-based storage for data persistence. Tests were conducted to verify:
\begin{itemize}
    \item Habit creation updates the storage file correctly
    \item A temporary test file is used during testing to avoid modifying real data
\end{itemize}
This ensures safe and reliable file handling.

---

\subsection{Task 7: Stability Testing}
The complete test suite was executed multiple times to verify stability. The results showed:
\begin{itemize}
    \item All tests passed consistently
    \item No random failures or inconsistencies occurred
\end{itemize}
This confirms that the system is stable and reliable.

---

\subsection{Task 8: Fault Detection Testing}
An intentional error was introduced into the validation logic of the create habit functionality. As a result:
\begin{itemize}
    \item The test suite failed and detected the issue
    \item A database error occurred due to missing required fields
\end{itemize}
After restoring the correct validation logic:
\begin{itemize}
    \item All tests passed successfully
\end{itemize}
This demonstrates that the test suite is effective in identifying bugs.

---

\subsection{Final Results}
The final test execution produced the following results:
\begin{itemize}
    \item Total Test Cases: 12
    \item Execution Status: All Passed
    \item System Status: Stable and Reliable
\end{itemize}

---


\section{Conclusion}

This project demonstrates how to build a complete REST-style API using Django 
with file-based persistence. The project helped in understanding backend 
development concepts including request handling, validation, persistence, 
and API design.

The final system is functional and can serve as a foundation for more advanced 
applications that use databases and user authentication.

\end{document}